\subsection{Datasets}
\label{sec:data-description}
The following is a more through description of the datasets used in this work.
We did not use any feature selection methods for any of the datasets. Although, for each dataset, 
we removed the uninformative features (e.g. various kinds of IDs). Categorical variables were converted 
to dummy/indicator variables, and an indicator for missing values were included per column with missing values.

\paragraph{The Adult dataset}
The Adult dataset~\citep{adult-dataset, uci-repo} from the UCI
repository contains Census data (hours worked per
week, education, sex, age, marital status, and so forth), and has been
used to train predictors as to whether an individual earns more or
less than \$50,000 per year. Predicting income based upon less
sensitive attributes, such as education and occupation, can be used to
aid in other economic decisions (related to lending or hiring). It is
therefore important to design models which are equally predictive of
income regardless of gender or race. We created groups for this dataset based
on gender. Furthermore, we extracted the data only from the Adult.data file.

\paragraph{The Communities and Crime dataset}
The Communities and Crime Dataset~\citep{communities-crime-dataset,
  uci-repo}, from the UCI repository is a dataset which includes many
features deemed relevant to violent crime rates (such as the
percentage of the community's population in an urban area, the
community's racial makeup, law enforcement involvement and racial
makeup of that law enforcement in a community, amount a community's
law enforcement allocated to drug units) for different
communities. This data is provided to train regression models based on
this data to predict the amount of violent crime (murder, rape,
robbery, and assault) in a given community. If police departments
allocate their units to communities based upon a prediction of this
form, this should draw concerns of fairness if the predictive behavior
is much better for violent predominantly white community than for a
violent predominantly nonwhite communities.  We consider communities
racial makeup as sensitive variable, however, the racial makeup
consists of more than two races.  We created two groups for this
dataset based on whether the percentage of black people in a community
(blackPerCap) is higher than the percentage of whites (whitePerCap),
indians (indianPerCap), asians (AsianPerCap) and hispanics
(HispPerCap) in that community.


%\paragraph{The student dataset}
%The Student Dataset~\citep{student-dataset}, from the UCI
%repository~\cite{uci-repo}, contains data relevant to secondary school
%achievement for students in Portugal, from information about students'
%home lives (their parents' education levels, jobs, cohabitation
%status, whether or not their homes have internet access, the length of
%their commute to school, etc), and educational information (whether
%additional paid courses were taken, the amount of time outside of
%school spent studying, whether nursery school was attended, etc). The
%goal from these attributes is to predict the grades of students in the
%final grading period in mathematics and Portugese, ideally without the
%use of the (highly correlated) grades from previous grading
%periods. If this prediction were used to allocate a limited resource
%to help struggling students, it would be important for that
%intervention to be allocated in similar ways across genders or races,
%conditioned on the individual students' likely grades. Similarly, if
%students were ``tracked'' based on these predictions (and placed in
%higher or lower achieving classrooms), it would be crucial that these
%predictions being similarly good in predicting different groups'
%performance.

\paragraph{The COMPAS dataset}
The COMPAS dataset contains data from Broward County, Florida originally compiled
by ProPublica~\cite{propublica} in which the goal
is to predict whether a convicted individual would commit a violent crime in the following two years
or not.
Following the analysis of Propublica and~\citet{Corbett-DaviesP17} we considered black
and white defendants who were assigned COMPAS risk scores
within 30 days of their arrest. Furthermore, we restricted ourselves to defendants who
``spent at least two years outside
a correctional facility without being arrested for a violent crime,
or were arrested for a violent crime within this two-year period~\cite{Corbett-DaviesP17}''.
Our goal is to predict the two-year violent recidivism. 
The dataset includes race, age, sex, number of prior convictions, and COMPAS violent crime
risk score (a score between 1 and 10 which we categorized to low, medium and high levels of risks).
As mentioned earlier, we created groups for this dataset based on race.

\paragraph{The Default dataset}
The Default of Credit Card Clients dataset~\citep{default-dataset, uci-repo} from the UCI
repository, contains data from Taiwanese credit card
users, such as their credit limit, gender, education, marital status,
history of payment, bill and payment amounts. These features are
designed to be used in the prediction of the probability of default
payments, although the data has binary labels. It has been
noted~\citep{default-dataset} that linear models tend to perform quite
poorly on this task. Given the sensitive nature of lending (and
historical inequity of the allocation of credit across races and
genders), it is important to understand the tradeoff between fairness
and accuracy in settings related to credit. We created groups for this 
dataset based on gender.


\paragraph{The Law School dataset}
The bar passage study was initiated in 1991 by Law School Admission Council national longitudinal. 
The dataset contains  records for law students who took the bar exam. 
The binary outcome indicates whether the student passed the bar exam or not. The features include variables such 
as cluster, lsat score, undergraduate GPA, zfyGPA, zGPA, full-time status, family income, age and also sensitive 
variables such as race and gender. 
The variable cluster is the result of a clustering of similar law schools (which is done apriori), and is used to adjust for the 
effect of type of law school. zGPA is the z-scores of the students overall GPA and zfyGPA is the  first year GPA relative to 
students at the same law school. We created groups for this dataset based on gender.
For more information about this dataset refer to \url{http://www2.law.ucla.edu/sander/Systemic/Data.htm}.

\paragraph{The Sentencing dataset}
The Sentencing dataset contains 5969 random sample of all prison inmates primarily from the department of corrections of a state. 
Moreover data from state agencies that have information of arrests are
also added to the dataset. To be sentenced to a state prison, an offender must have been convicted of a felony or must have pled guild to a felony. 
The outcome variable, sentence length, is the sentence given by a judge, but is not necessarily the amount of time served in prison. 
There are usually provisions for review by the state parole board, which can lead to release well before the full nominal sentence is served. 
The predictors include variables that are legitimate factors in sentencing (e.g., the crime for which a conviction was obtained) as well as sensitive variables 
like gender.  We created groups for this dataset based on gender.






